\scanpage{21}
\diagram{\begin{tikzcd}[ampersand replacement=\&,column sep=tiny,row sep=small]
 A \arrow[dr,dash] \& \& B \arrow[dl,dash]\\
 \& S \&
\end{tikzcd}}
\margin{schémas abéliens sur $S$}
\[
u_p:T_p(A)\to T_p(B)
\]
Question: quand cela provient-il d'un hom $u:A\to B$.

\deleted{(i) A cause de l'unicité, la descente fid. plate marche.

(ii) l'ens \inserted{$T$} des pts $s\in S$ tels que $u_{ps}:T_p(A_s)\to T_p(B_s)$ provienne de $u_s:A_s\to B_s$ (condition géom. par (i)!) est \uline{fermé}.

\deleted{Si $t\in T$ est tel que $\operatorname{car.} k(t)=$}

(iii) Si $t\in T$ est tel que $\operatorname{car} k(t)=p$, alors $\exists$ voisinage ouvert $U$ de $t$ tel que $u_p\mid U$ soit algébrique, en particulier \deleted{\unsure{Tout|$t$ est|$T$}} \unsure{$t$ y est|y est|$U$ est} pt intérieur de $T$.}

\diagram{\begin{tikzcd}[ampersand replacement=\&,column sep=small,row sep=small,execute at end picture={
\draw (-2.9,-1.75) .. controls (-2.2,.10) and (-1.4,1.20) .. (-.9,1.65);
\draw (-2.25,-1.72) .. controls (-1.75,-.3) and (-1.0,.9) .. (.15,1.65);
\draw (-1.5,-1.62) .. controls (-1.1,-.1) and (.15,1.2) .. (1.1,1.65);
\draw (-.83,-1.50) .. controls (-.4,-.2) and (1.1,1.15) .. (2.2,1.57);
\draw (-2.9,1.0) -- (-2.9,-1.75) .. controls (-1,-1.7) and (-.12,-1.5) .. (.1,-1.4) .. controls (.6,-1.2) and (.15,-.6) .. (.1,-.4) .. controls (.8,-.6) and (1.4,-.56) .. (2.1,-.55) -- (2.25,.92);
\draw (-2.62,-1.79) circle (.18); \fill (-2.62,-1.79) circle (.035); \draw (-3.15,-1.79) -- (-2.62,-1.79);
}]
 \& \& \& \deleted{A_{\unsure{s'|S'}}} \arrow[r] \& \deleted{B_{s'}}\\
 \deleted{S} \& \deleted{\unsure{s|S_0}} \& \deleted{\unsure{S''|S'''|S_1}} \& \deleted{S'} \arrow[l] \&
\end{tikzcd}}

a) Il y a \inserted{(si $S$ noeth)} un plus grand sous-schéma $T$ de $S$, sur lequel $u_p$ est algébrique ; $T$ est fermé, sa formation commute à tout changt de base (nouvelles bases \unsure{pas néc.|par ex.|par voisinages} noethériennes).

b) Si $t\in T$ est tel que $\operatorname{car} k(t)=p$, alors $T$ majore un voisinage ouvert induit de $t$. \deleted{Donc si $S_p=V(p1_S)$, alors $T$ est un voisinage \unsure{(schématique)} de $T\cap S_p$}

\deleted{\uline{Cor.} Supposons que pour tout $s\in\operatorname{Ass}\mathcal O_S$,}

\deleted{\uline{Corollaire} Supposons que pour tout $s\in$}

\uline{Cor.} Soit $Z$ l'ens \deleted{\unsure{des réunion} de $\operatorname{Max} S$ et de (pts maximaux de $S$) et de l'ens des} \inserted{des pts de $S$ qui sont} maximaux, ou qui sont associés \deleted{\unsure{et}} à $\mathcal O_S$ \uline{et} de car. $\ne p$. Supposons que $z\in Z\Longrightarrow\exists\,s\in\overline{\{z\}}$ tel que $k(s)$ de car. $p$ et $s\in T$ \inserted{(i.e. $\overline{\{z\}}\cap S_p\cap T\ne\varnothing$)}. Alors $u_p$ est algébrique.
% archive sheet number: 20

